\documentclass[10pt,a4paper]{article}
\usepackage[margin=0.79in]{geometry}
\usepackage{amsmath,amssymb,booktabs}
\usepackage[hidelinks]{hyperref}
\setlength{\parindent}{0pt}
\setlength{\parskip}{5pt}
\newcommand{\head}[1]{\par\medskip\textbf{#1}\par}
\begin{document}
\begin{center}
{\large\bfseries Congestion budgets at the selected lane allocator}\\[3pt]
2026-10-21 --- conditional interface check; LF2 unselected
\end{center}

\textbf{Scope.} The user selects the attempt-and-forward allocator (LF1).
This note checks its demand/receiving inputs against the Subraveti summary
and cited Han formulation. It isolates a consequential sending-mode choice;
it does not complete the task, select LF2 or run a traffic trajectory.
The formulas below cover a common triangular fundamental diagram (FD), not
the source's heterogeneous-lane state mapping.

\head{1. The complete source-lineage budget on a common FD}
Let nominal free speed $u_0>0$, backward wave-speed magnitude $w>0$, and jam
density $k_J>0$ be fixed geometry/traffic parameters. Let $0\leq\delta<1$ be
the maximum capacity-drop fraction. This $\delta$ is the source's capacity-drop
parameter, not Carey's lane-change headway multiplier. Define
\begin{align}
 C&=\frac{u_0wk_J}{u_0+w}, & k_c&=C/u_0,\nonumber\\
 Q_J&=(1-\delta)C, & k_{c,J}&=Q_J/u_0, &
 w_d&=\frac{Q_J}{k_J-k_{c,J}}
      =\frac{u_0w(1-\delta)}{u_0+\delta w}. \label{eq:fixed}
\end{align}
$w_d$ is the positive recovery-branch slope, with $0<w_d\leq w$.
The primed critical density in Han 2016's coefficient is $k_{c,J}$, evaluated
at maximum jam/drop; it is not a new advice-dependent parameter
\cite[sec. 3.2, eqs. (3)--(7)]{han16}. Han 2017 supplies the explicitly
capped ceiling and the VSL interpretation \cite[sec. 2.3]{han17}:
\begin{equation}
 Q(x)=\min\left\{C,\ C\left[1-\delta\frac{x-k_c}{k_J-k_c}\right]\right\},
 \qquad 0\leq x\leq k_J. \label{eq:ceiling}
\end{equation}
Consequently $Q_J\leq Q(x)\leq C$. Its input $x$ is the density of the
preceding same-lane cell. The min with $C$ follows the source's no-drop
under-critical branch; it is not an unapproved numerical clipping repair.

For consecutive same-lane densities $x=k_{i,l}$ and $y=k_{i+1,l}$, the
common-FD receiving law is
\begin{equation}
 R_H(x,y)=\min\{Q(x),\ w(k_J-y),\ w(k_J-x)+w_d(x-y)\}. \label{eq:receive}
\end{equation}
This is the original Han min-of-three form with its constants resolved,
consistent with Han 2017's piecewise supply. It uses both adjacent upstream
and downstream traffic, not only the downstream cell's empty space.

\textbf{H: the full-lineage sending law.} For advice $v_{i,l}$,
\begin{equation}
 D^H_{i,l}(v)=\min\{v_{i,l}k_{i,l},\ Q(k_{i-1,l})\},\qquad
 0<v_{\min}\leq v_{i,l}\leq u_0. \label{eq:demandH}
\end{equation}
The source uses $Q_{i,l}=Q(k_{i-1,l})$: its index identifies the cell whose
outflow is limited; its pressure input comes from the preceding cell. VSL
restricts sending and does not replace $u_0$ in \eqref{eq:fixed}, move the
physical grid or directly alter receiving/priority weights. A fixed control
mask can apply common advice in the designated upstream cells and use
$v=u_0$ elsewhere. Its exact task boundary still needs consolidation.

\head{2. The shared-receiver interpretation}
For movements into $(i+1,A)$, use \emph{one} budget
$\Delta t\,R_H(k_{i,A},k_{i+1,A})$ and share it among all allowed donors.
Do not give each donor a separate copy. This retains the principal source's
lane index and the cited source's shared-mainline-supply structure. E's
congestion does not independently become a second direct pressure argument.
The selected endpoint sends last-E mass only toward A; its donor ceiling
$Q(k_{i-1,E})$ remains defined without a downstream E cell. Boundary ghost
states for the first physical row remain a task-boundary obligation.

\newpage
\head{3. Why the alternative is a model choice}
The principal paper~\cite[eqs. (18)--(19) and following paragraph]{ns}
says that modified demand and supply apply in a discharging state
$k_{i,l}\geq k_{i+1,l}$, and ordinary functions otherwise. The thesis repeats
the sentence~\cite[chap. 3, eqs. (3.18)--(3.19)]{thesis}. One literal joint
reading yields the following \emph{G completion}:
\begin{equation}
 D^G_{i,l}(v)=
 \begin{cases}
 \min\{v_{i,l}k_{i,l},Q(k_{i-1,l})\},&k_{i,l}\geq k_{i+1,l},\\
 \min\{v_{i,l}k_{i,l},C\},&k_{i,l}<k_{i+1,l}.
 \end{cases} \label{eq:demandG}
\end{equation}
It completes the otherwise abbreviated constants from the original source;
it is not a claim about unpublished author code. To complete G at the last
E cell, whose same-lane successor does not exist, the proposed convention
keeps the first branch enabled for its compulsory merge. That endpoint
convention is part of candidate G, not an approved default.

Use \eqref{eq:receive} for the common receiving budget in both candidates.
On this common FD, when $x\leq y$, it reduces to
$\min\{C,w(k_J-y)\}$; the source's non-discharging receiving branch is thus
recovered. The sending switch in \eqref{eq:demandG} is the consequential
difference isolated here. Both candidates can conserve mass. The source's
short summary and full ancestor are insufficient to identify author intent
without selecting or clarifying this difference explicitly.

\head{4. Exact static case: a receiving queue grows, but G serves more}
Use the earlier illustrative FD, not a calibration:
\begin{equation}
 u_0=60\ {\rm km/h},\quad w=20\ {\rm km/h},\quad k_J=150\ {\rm veh/km},
 \quad C=2250\ {\rm veh/h},\quad k_c=37.5,\quad\delta=0.1.
\end{equation}
Then $Q_J=2025$, $k_{c,J}=33.75$, $w_d=540/31$ km/h. Set
\begin{equation}
 k_{i-1}=150,\quad k_i=40,\quad v_i=60,\qquad
 k_{i+1}=40-\varepsilon\ \hbox{or}\ 40+\varepsilon.
\end{equation}
All densities are in veh/km. No lateral requests are used: this is the
one-lane reference reduction of the allocator. We have $Q(150)=2025$,
$Q(40)=2245$ and $v_i k_i=2400$. For small positive $\varepsilon$,
\begin{align}
 R_H(40,40-\varepsilon)&=2200+\frac{540}{31}\varepsilon>2025,\nonumber\\
 R_H(40,40+\varepsilon)&=2200-20\varepsilon>2025. \label{eq:receivesides}
\end{align}
For example these expressions apply at $\varepsilon=1/10$. Hence
\begin{center}
\begin{tabular}{lrr}
\toprule
 & $k_{i+1}=39.9$ & $k_{i+1}=40.1$\\
\midrule
$D^G_i$ (veh/h) & 2025 & 2250\\
$\min\{D^G_i,R_H\}$ (veh/h) & 2025 & 2198\\
$\min\{D^H_i,R_H\}$ (veh/h) & 2025 & 2025\\
\bottomrule
\end{tabular}
\end{center}
As $\varepsilon\downarrow0$, G's flow limits differ by $2200-2025=175$
veh/h; H's do not. The difference is not hidden by receiving. The example
does not depend on evaluating a missing supply branch away from its domain:
the recovery-side formula is supplied by the original Han source, and the
buildup side uses the ordinary space bound. It diagnoses a literal sending
interpretation, not the authors' realized implementation.

\newpage
\head{5. Conditional budget and storage qualification}
For $0\leq x,y\leq k_J$, rewrite the third receiving term as
\begin{equation}
 w(k_J-x)+w_d(x-y)
 =w_d(k_J-y)+(w-w_d)(k_J-x)\geq0. \label{eq:positive}
\end{equation}
Together with \eqref{eq:ceiling}, this proves
\begin{equation}
 0\leq R_H(x,y)\leq w(k_J-y),\qquad
 0\leq D^H_{i,l},D^G_{i,l}\leq v_{i,l}k_{i,l}. \label{eq:bounds}
\end{equation}
On fixed equal-length cells, choose
\begin{equation}
 u_0\Delta t\leq\Delta x,\qquad w\Delta t\leq\Delta x.
\end{equation}
Let $n_{i,l}=\Delta x k_{i,l}$ and $J=\Delta x k_J$. Then the allocator
inputs $d_{i,l}=\Delta t D_{i,l}$ and $r_{i,l}=\Delta t R_{i,l}$ obey
\begin{equation}
 0\leq d_{i,l}\leq n_{i,l},\qquad
 0\leq r_{i,l}\leq J-n_{i,l}.
\end{equation}
The selected allocator's common donor and receiver bounds therefore imply
nonnegative stocks, jam-storage protection and conservation for either
candidate. Its one-cell diagonal connections and no-instant-reuse update
are unchanged. Setting $\delta=0$ gives $w_d=w$, $Q=C$ and ordinary CTM
budgets, so G and H coincide in the no-capacity-drop limit.

\textbf{A helpful source caution.} Subraveti motivates its same-cell scheme
by warning that diagonal transfer depending only on downstream supply can
misrepresent congestion propagation. The selected geometry changes that
scheme, but \eqref{eq:receive} and the retained lookahead incentives still
use adjacent/upstream lane information. This addresses information retained
in the join; it does \emph{not} prove equivalent wave propagation or inherit
the source's calibration. No traffic run tests that claim here.

\head{6. Sensitivity scope and decision}
H's budgets are continuous, piecewise smooth functions of their current
density inputs and advice. With fixed active minima and request-allocation
sets they can be differentiated normally. G adds a possible finite jump at
the density-order guard; a derivative through just the selected arithmetic
branch does not capture crossing that jump. Neither statement validates a
whole rollout gradient. The native incentives, empty states, time-boundary
choices and complete population-time objective remain to be reviewed.

\textbf{LF2 recommendation: H.} It uses the full congestion law in the
principal model's own source lineage and the already proposed Han VSL map.
It avoids the extra sending switch and endpoint convention. Selecting it
still resolves a scientific interpretation affecting all movements; it is
not silently approved as a transcription fix. LF1 stays selected. The
review stops here before full task/information/accounting/gradient closure.

\begin{thebibliography}{9}\setlength{\itemsep}{0pt}
\bibitem{ns} H. H. S. Nagalur Subraveti, V. L. Knoop, and B. van Arem,
``First order multi-lane traffic flow model -- an incentive based macroscopic
model to represent lane change dynamics,'' \emph{Transportmetrica B},
vol. 7, no. 1, pp. 1758--1779, doi: 10.1080/21680566.2019.1700846.
\bibitem{han16} Y. Han, Y. Yuan, A. Hegyi, and S. P. Hoogendoorn,
``New extended discrete first-order model to reproduce propagation of jam
waves,'' \emph{Transp. Res. Rec.}, vol. 2560, pp. 108--118, 2016,
doi: 10.3141/2560-12. Institutional accepted manuscript, sec. 3.2.
\bibitem{han17} Y. Han \emph{et al.}, ``Resolving freeway jam waves by
discrete first-order model-based predictive control of variable speed limits,''
\emph{Transp. Res. C}, vol. 77, pp. 405--420, 2017,
doi: 10.1016/j.trc.2017.02.009. Accepted manuscript, sec. 2.3.
\bibitem{thesis} H. H. S. Nagalur Subraveti, \emph{Lane-specific Traffic
Flow Control}, Ph.D. dissertation, TU Delft, 2021,
doi: 10.4233/uuid:e405a549-b1fe-4fa1-8090-4429612359d9.
\end{thebibliography}
\end{document}
